\section{Introduction to AI-Generated Text Forensics}
\noindent The rapid evolution of generative artificial intelligence has led to a significant challenge in digital forensics: distinguishing human-written text from machine-generated content. As Large Language Models (LLMs) grow in size and complexity, they generate text that exhibits high semantic coherence, correct grammar, and contextual relevance. While this facilitates automation in content creation and programming, it introduces substantial risks in academic integrity, information security, and public trust. 

Text forensics has emerged as a crucial field to mitigate these risks. Current methodologies can be categorized into three primary tasks:
\begin{enumerate}
    \item \textbf{Binary Detection:} Distinguishing between human-written and machine-generated text.
    \item \textbf{Closed-Set Attribution:} Classifying machine-generated text into one of several predefined generator models (e.g., attributing a text specifically to GPT-4 or LLaMA-2).
    \item \textbf{Open-Set Generator Rejection:} Identifying whether a text was generated by a model that was completely unseen during training, and rejecting it as an unknown source rather than misclassifying it into a known class.
\end{enumerate}

\section{Supervised and Lexical Detection Paradigms}
\subsection{Lexical and N-Gram Methods}
Early text forensics focused on statistical word and character distributions. Lexical classifiers, such as FastText \cite{joulin2016bag}, utilize bag-of-words representations and n-gram frequencies to fit simple linear classifiers. While these models are highly efficient and scale to massive datasets with minimal computational cost, they suffer from key limitations:
\begin{itemize}
    \item \textbf{Contextual Blindness:} They treat text as unordered bags of tokens, ignoring syntax, word ordering, and semantic coherence.
    \item \textbf{Vulnerability to Synonyms:} Simple lexical substitutions (e.g., using synonyms) can easily bypass n-gram detectors.
\end{itemize}

\subsection{Fine-Tuned Transformer-Based Classifiers}
To capture context, researchers fine-tune pre-trained bidirectional transformer encoders (such as BERT \cite{devlin2018bert} or DistilBERT \cite{sanh2019distilbert}) on balanced corpora of human and machine texts. In these models, the representation of the special classification token (\texttt{[CLS]}) in the final layer aggregates the global context of the entire sequence through self-attention.
Supervised transformer models achieve state-of-the-art accuracy on closed-set benchmarks. However, they are highly susceptible to out-of-distribution (OOD) shifts:
\begin{itemize}
    \item \textbf{Domain Shifts:} If a model is trained on news summaries and evaluated on creative writing, its accuracy drops significantly.
    \item \textbf{Generator Shifts:} When exposed to text generated by a model family absent from the training set, supervised classifiers fail to generalize, leading to high false attribution rates.
\end{itemize}

\section{Zero-Shot and Predictability-Based Detectors}
\subsection{Token Probability Predictability}
Zero-shot methods evaluate text without training classifiers on labeled data, relying instead on the internal token probability distributions of a reference language model (like GPT-2). GLTR (Giant Language Model Test Room) \cite{gehrmann2019gltr} represents the foundational work in this category. It computes the rank of each token under the reference model's prediction distribution. Because language models select highly predictable words (top ranks) to maintain fluency, machine-generated text exhibits lower rank entropy compared to human writing, which contains more unpredictable word choices (tail ranks).

\subsection{Curvature-Based Perturbation Methods}
DetectGPT \cite{mitchell2023detectgpt} improves upon probability ranks by evaluating the local curvature of the reference model's log-likelihood space. It perturbs the candidate text by replacing random words with synonyms using a mask-filling model (e.g., DistilRoBERTa) and measures the resulting drop in log-likelihood under the reference LLM. If the original text lies in a local peak of the model's likelihood space, the perturbed texts will show a sharp decrease in likelihood, signaling machine generation. While mathematically robust and require no training, perturbation methods are highly computationally expensive, requiring hundreds of model evaluations per text.

\section{Open-Set Out-of-Distribution (OOD) Rejection}
\subsection{The Open-Set Challenge in Attribution}
Closed-set classifiers operate under the assumption that all evaluation samples belong to the classes present in the training set. In practical scenarios, however, new generator models are continually released. When a closed-set classifier encounters text from an unseen model (e.g., OPT), it is forced to classify it as one of the known models (e.g., LLaMA), causing severe attribution errors.

\subsection{Distance-Based OOD Detection}
To handle open-set attribution, researchers map text representations into high-dimensional embedding spaces (like DistilBERT's 768-D space) and model the distributions of the known classes. The out-of-distribution score is computed based on the distance from the input embedding to the nearest known class distribution.
Traditional Euclidean distance fails in high-dimensional spaces because it treats all embedding dimensions as independent and isotropic. Angular Cosine distance measures only the direction of the vectors, ignoring magnitude variations.
Mahalanobis distance \cite{mahalanobis1936generalized} resolves these issues by scaling the distance by the covariance matrix of the known training distributions:
$$D_M(z) = \min_{c} \sqrt{(z - \mu_c)^T \Sigma_c^{-1} (z - \mu_c)}$$
where $\mu_c$ is the class centroid and $\Sigma_c$ is the regularized class covariance. The covariance scaling accounts for multi-dimensional correlations and variances, expanding the space along direction of high class variation.

\enlargethispage{2\baselineskip}
\subsection{Threshold Selection Strategies}
Selecting an OOD threshold $\tau$ requires balancing the rejection of unseen data against the false rejection of known data. The standard strategy, introduced by Hendrycks and Gimpel \cite{hendrycks2016baseline}, selects $\tau$ as the 95th percentile of the minimum Mahalanobis distances on a known validation set. This guarantees that at least 95\% of the known samples are correctly classified, while establishing a strict boundary to reject unseen generator families as \textbf{UNKNOWN}.
